\section{Conclusion}
\label{sec:conclusion}

This paper has presented PC-Diffusion, a novel physics-constrained generative framework targeting Zero-Shot Fault Diagnosis (ZSFD) in rotating machinery, actively answering the demand for reliable and interpretable industrial diagnostics. Recognizing the limitations of standard unconstrained generative models, our methodology embeds the mathematical fundamentals of mechanical vibrations directly into the diffusion training paradigm. Synchronously, the introduction of a dual-level Bayesian semantic embedding fundamentally resolves the previously elusive issue of diagnostic uncertainty quantification in test-time scenarios. Validated across empirical benchmarks, the architecture establishes a new standard for generation quality and zero-shot diagnostic resilience. 

Future research will focus on several promising directions. First, we aim to optimize the computational efficiency of PC-Diffusion for real-time deployment on edge devices with limited resources, enabling ultra-low latency fault detection in industrial Internet of Things (IIoT) environments. Second, we plan to extend the framework to incorporate multi-sensor fusion, integrating vibration, acoustic, and thermal data streams to enhance diagnostic robustness under complex operating conditions. Finally, exploring the application of PC-Diffusion to compound fault scenarios and multi-component systems represents a critical step toward comprehensive industrial health monitoring.